\documentclass[12pt,a4paper]{article}

% =====================================================
% PACOTES E CONFIGURAÇÕES DE PÁGINA
% =====================================================
\usepackage[utf8]{inputenc}
\usepackage[T1]{fontenc}
\usepackage[brazil,shorthands=off]{babel}
\usepackage{geometry}
\usepackage{graphicx}
\usepackage{fancyhdr}
\usepackage{titlesec}
\usepackage[table]{xcolor}
\usepackage{hyperref}
\usepackage{setspace}
\usepackage{enumitem}
\usepackage{tabularx}
\usepackage{booktabs}
\usepackage{newtxtext,newtxmath}
\usepackage[most]{tcolorbox}
\usepackage{listings}
\usepackage{float}

\geometry{
    top=2.2cm,
    bottom=2.2cm,
    left=2.0cm,
    right=2.0cm
}

\onehalfspacing

% =====================================================
% CORES INSTITUCIONAIS (IFSC)
% =====================================================
\definecolor{ifscGreen}{RGB}{22,163,74}
\definecolor{azulTitulo}{RGB}{30,60,90}
\definecolor{cinzaCabecalho}{RGB}{240,244,248}
\definecolor{cinzaCodigo}{RGB}{248,250,252}
\definecolor{borderGray}{RGB}{203,213,225}

% =====================================================
% CONFIGURAÇÃO DE LISTINGS (CÓDIGO LATEX)
% =====================================================
\lstset{
    backgroundcolor=\color{cinzaCodigo},
    basicstyle=\ttfamily\small,
    breaklines=true,
    frame=single,
    rulecolor=\color{borderGray},
    keywordstyle=\color{azulTitulo}\bfseries,
    commentstyle=\color{gray}\itshape,
    stringstyle=\color{ifscGreen},
    numbers=none,
    showstringspaces=false,
    tabsize=2
}

% =====================================================
% CONFIGURAÇÃO DE HIPERLINKS
% =====================================================
\hypersetup{
    colorlinks=true,
    linkcolor=azulTitulo,
    urlcolor=ifscGreen,
    citecolor=azulTitulo
}

% =====================================================
% ESTILOS DE SEÇÃO
% =====================================================
\titleformat{\section}
{\normalfont\large\bfseries\color{azulTitulo}}
{\thesection}
{1em}
{}
[\vspace{0.15cm}\hrule\vspace{0.2cm}]

\titleformat{\subsection}
{\normalfont\normalsize\bfseries\color{azulTitulo}}
{\thesubsection}
{1em}
{}

% =====================================================
% CABEÇALHO E RODAPÉ
% =====================================================
\pagestyle{fancy}
\fancyhf{}
\fancyhead[L]{\small\bfseries\color{azulTitulo} IFSC Câmpus Garopaba • CST Sistemas para a Internet}
\fancyhead[R]{\small\color{gray} Iniciação Científica \& Extensionista}
\fancyfoot[C]{\thepage}
\fancyfoot[L]{\footnotesize\color{gray} Roteiro 1 — Overleaf, ACM COTB, Imagens e Tabelas}
\fancyfoot[R]{\footnotesize\color{gray} Semestre 2026/2}
\renewcommand{\headrulewidth}{0.5pt}
\renewcommand{\footrulewidth}{0.5pt}

\begin{document}

% =====================================================
% CABEÇALHO DO DOCUMENTO
% =====================================================
\begin{tcolorbox}[
    colback=cinzaCabecalho,
    colframe=borderGray,
    arc=3mm,
    boxrule=1pt,
    left=15pt,
    right=15pt,
    top=12pt,
    bottom=12pt
]
\begin{center}
    {\large\bfseries\color{azulTitulo} INSTITUTO FEDERAL DE SANTA CATARINA (IFSC)}\\
    {\small CÂMPUS GAROPABA — DEPARTAMENTO ACADÊMICO}\\
    \vspace{0.15cm}
    {\bfseries CURSO SUPERIOR DE TECNOLOGIA EM SISTEMAS PARA A INTERNET}\\
    \vspace{0.25cm}
    {\Large\bfseries\color{ifscGreen} ROTEIRO DE PRÁTICA EM LABORATÓRIO Nº 01}\\
    \vspace{0.15cm}
    {\large\bfseries Estruturação de Artigos no Padrão ACM / COTB:\\Manipulação de Imagens, Tabelas e Referências Cruzadas no Overleaf}
\end{center}
\vspace{0.2cm}
\noindent\textbf{Unidade Curricular:} Iniciação Científica e Extensionista (60h — 30h Extensão)\\
\textbf{Docentes Responsáveis:} Prof. André Moraes \& Prof. Adalberto Tabalipa\\
\textbf{Meta Final da UC:} Artigo Científico para o \textit{Computer on the Beach} (COTB)
\end{tcolorbox}

\vspace{0.3cm}

\section{Apresentação e Contextualização}

O evento científico \textbf{Computer on the Beach (COTB)} adota o modelo internacional de publicação da \textbf{ACM (Association for Computing Machinery)}, baseado na classe \texttt{acmart} com diagramação em \textbf{duas colunas}. 

Diferente de editores convencionais (como o MS Word ou Google Docs), no \LaTeX{} o autor não arrasta elementos visualmente. O compilador calcula a melhor posição tipográfica baseado em regras estritas de paginação. Portanto, dominar o controle de \textbf{largura relativa} e os ambientes de \textbf{1 coluna versus 2 colunas} é indispensável para a produção do artigo final da disciplina.

\section{O Template Oficial no Overleaf}

Para iniciar o trabalho, cada dupla deve acessar o template oficial da ACM no Overleaf:
\begin{center}
    \url{https://www.overleaf.com/latex/templates/association-for-computing-machinery-acm-sig-proceedings-template/bmvfhcdnx0px}
\end{center}

\noindent O arquivo principal (\texttt{main.tex}) deve iniciar com a declaração:
\begin{lstlisting}[language=TeX]
\documentclass[sigconf,nonacm]{acmart} % Modo conferência em duas colunas
\end{lstlisting}

\section{Manipulação de Imagens (Figuras)}

No modelo em duas colunas, existem duas formas de incluir figuras:

\subsection{1. Imagem na Largura de Uma Coluna (\texttt{figure})}
Usada para gráficos, capturas de tela verticais e ilustrações pontuais. Utiliza-se a medida \texttt{\textbackslash columnwidth}:

\begin{lstlisting}[language=TeX]
\begin{figure}[htbp]
  \centering
  \includegraphics[width=\columnwidth]{grafico_resultados.png}
  \caption{Média do tempo de resposta obtida nos testes experimentais.}
  \label{fig:tempo_resposta}
\end{figure}
\end{lstlisting}

\subsection{2. Imagem Panorâmica / Largura Total da Página (\texttt{figure*})}
Quando o diagrama da arquitetura ou o fluxo metodológico for muito detalhado, utiliza-se o ambiente \texttt{figure*} (com asterisco) e a medida \texttt{\textbackslash textwidth}:

\begin{lstlisting}[language=TeX]
\begin{figure*}[htbp]
  \centering
  \includegraphics[width=0.92\textwidth]{arquitetura_pipeline.png}
  \caption{Visão geral do pipeline da aplicação e fluxo de interação com a comunidade.}
  \label{fig:arquitetura_geral}
\end{figure*}
\end{lstlisting}

\begin{tcolorbox}[colback=yellow!10,colframe=orange!80!black,arc=2mm,title=\textbf{Regra de Ouro da Redação Científica:}]
\small Nunca utilize expressões como \textit{"na figura abaixo"} ou \textit{"na tabela acima"}. Sempre utilize o comando de referência cruzada com til (espaço inquebrável):\\
\texttt{Conforme apresentado na Figura\textasciitilde\textbackslash ref\{fig:tempo\_resposta\}...}
\end{tcolorbox}

\section{Manipulação de Tabelas Acadêmicas}

O modelo da ACM desaconselha linhas verticais pesadas e incentiva o uso do pacote \texttt{booktabs} com linhas horizontais limpas (\texttt{\textbackslash toprule}, \texttt{\textbackslash midrule}, \texttt{\textbackslash bottomrule}).

\subsection{1. Tabela em 1 Coluna}
\begin{lstlisting}[language=TeX]
\begin{table}[htbp]
  \caption{Comparativo das ferramentas analisadas na pesquisa}
  \label{tab:comparativo_ferramentas}
  \centering
  \begin{tabular}{lcr}
    \toprule
    \textbf{Ferramenta} & \textbf{Tipo de Coleta} & \textbf{Tempo Médio (s)} \\
    \midrule
    Lighthouse API      & Automatizada           & 3,45 \\
    Script Python       & API / Requisições      & 1,20 \\
    Auditoria Manual    & Heurística             & 15,80 \\
    \bottomrule
  \end{tabular}
\end{table}
\end{lstlisting}

\subsection{2. Tabela Panorâmica (2 Colunas com \texttt{table*})}
Para quadros comparativos extensos de Trabalhos Relacionados:
\begin{lstlisting}[language=TeX]
\begin{table*}[htbp]
  \caption{Mapeamento Sistemático da Literatura e Estado da Arte}
  \label{tab:estado_da_arte}
  \centering
  \begin{tabular}{p{3.5cm} p{4.5cm} p{4.5cm} c}
    \toprule
    \textbf{Referência} & \textbf{Objetivo do Estudo} & \textbf{Limitações Identificadas} & \textbf{Ano} \\
    \midrule
    Silva e Moraes \cite{silva2024} & Análise de acessibilidade em portais & Amostra restrita à capital & 2024 \\
    Blankenburg et al. \cite{blankenburg2025} & Ferramenta com Raspberry Pi & Foco exclusivo no público infantil & 2025 \\
    \bottomrule
  \end{tabular}
\end{table*}
\end{lstlisting}

\section{Atividade Prática em Laboratório (Entregáveis de Hoje)}

Cada dupla (ou estudante individual) deverá realizar no Overleaf:
\begin{enumerate}[leftmargin=1.8cm, label=\textbf{Passo \arabic*:}]
    \item Criar o projeto no Overleaf a partir do template ACM (\texttt{sigconf}).
    \item Configurar o título do artigo da dupla, nome dos autores e afiliação institucional (\textit{Instituto Federal de Santa Catarina}).
    \item Fazer o upload de uma imagem ou gráfico e inseri-la usando \texttt{\textbackslash includegraphics}, \texttt{\textbackslash caption} e \texttt{\textbackslash label}.
    \item Construir uma tabela com dados do tema da dupla formatada com o pacote \texttt{booktabs}.
    \item Citar a imagem e a tabela no texto corrido da Introdução/Metodologia usando \texttt{\textbackslash ref\{...\}}.
    \item Obter o \textbf{link de compartilhamento em modo leitura (\texttt{/read/})} e atualizar no Painel de Equipes do Dashboard da UC.
\end{enumerate}

\end{document}
